% \citep{basu-raga}
% * Two stage deep learning approach
% * MFCC

% \citep{jayanthi-raga}
% * Features: pre-processing of audio emphasizing spectral centroid, and then next, spectral bandwidth, spectral roll-off, zero crossing rate, and Mel-Frequency Cepstral Coefficients (MFCC) are used to extract features.
% * BiLSTM-XGBoost
% * no pitch no separation

% \citep{sridhar-raga}
% * svara table per artists, frequency features from harmonic components of the spectrogram
% * separation using \citep{zhang-separation}, under the conclusions of \citep{sridhar-comparison}

% \citep{koduri-survey}
% * rao pitch extraction, no separation
% * note segmentation + pitch class profile (histogram)

% \citep{shetty-raga}
% * scale matching
% * Autocorrelation for pitch. For simplicity, the type of the audio file taken as input is monophonic song with a single sound-source. System can be enhanced further for recorded Polyphonic music, by applying multi-pitch detection techniques described in [5]. 

% \citep{shah-raga_recognition}
% * spleeter, not clear improvement tbh, not pitch but spectrograms
% * CNN + LSTM
% * 10 Hindustani ragas


% \citep{kumar-raga}
% * melodia
% * no separation
% * We evaluated our proposed approach on CompMusic dataset and our own dataset and show that combining the clues from pitch-class profiles and n-gram histogram indeed improves the performance (use SVM to classify in ragas).

% \citep{gulati-phrase_raga}
% * uses same input as \citep{gulati-mining} --> (melodia, processing, no separation)
% * removing unvoiced segments using classification models

% \citep{anand-raga}
% * no separation, 
% * melodia
% * CNN
% * Chromogram, learning features from predominant pitch values from mix
% * 16 Carnatic ragas

% \citep{chowdhuri-raga}
% * chromagrams, not pitch directly, deep learning (CNN, and RNN to get context)
% * 30 Hindustani raagas

% \citep{belle-raga}
% * Pitches were extracted every 20 ms from the range of 100 Hz to 1000 Hz with a resolution of 0.01 Hz. The obtained pitch contour was validated by listening to the re-synthesised pitch contour. Any vocal detection errors or pitch tracking errors were corrected by selecting the specific segments of the input audio and running the melody extractor with manually adjusted parameters. Accurate pitch contours corresponding to the vocal melody were thus extracted for all the segments in the study.
% * Rao method
% * pitch class distribution
% * 4 raagas

% \citep{gulati-time_delayed}
% * melodia, no separation
% * time delayed melody surfaces (TDMS)
% * 30 Hindustani ragas, 40 Carnatic ragas

% \citep{pandey-tansen}
% * based on Hidden Markov Model and string, Pakad matching. Their test data consists of result on only 2 ragas Yaman kalian and Bhupali. They used HMM model because notes are small in number and the sequence for raga is very well defined
% * Note transcription is done through two heuristics which convert input audio sample into sequence of notes. The two heuristics are as follows: (i) The Hill Peak Heuristic: This heuristic identifies notes from input audio sample based on hill and peak value in its pitch graph. (ii) The Note Duration Heuristic: First of all the assumption is made that a composition of music continues for at least some constant amount of time or duration
% * no pitch, no gamakas, no separation, only 2 ragas

% \citep{koduri-raga}
% * Combination of YIN + Melodia. But this has an inherent quantization step which does not allow high bin resolutions in histogram analysis. So we use a combination of both. In each frame, we transform the estimated pitch value from both methods into one octave and compare them. In those frames where they agree within a threshold, we retain the corresponding YIN pitch transforming it to the octave of the pitch-value from multi-pitch analysis. We discard the data from frames where they disagree with each other.
% * No separation
% * Histograms basically

% \citep{chordia-raga}
% * Then pitch detection is done on sample segment using Harmonic Product Spectrum (HPS) algorithm. Thresholding complex Detection Function (DF) on each segment, note onsets are determined. 
% * Pitch class distribution
% * 31 Hindustani ragas (13 vocalists, 6 instruments)

% \citep{shikarpur-raga_classification}
% * Analysis of mode switching in Hindustani music, the first one. Use pitch-based features
% * Use the four-stem spleeter
% * Use the Python API of PRAAT

% \citep{ross-raga_similarity}
% * note embededings from a dataset: Our experiments are carried out with the Hindustani bandish dataset available from swarganga.org, created by Swarganga music foundation. This website is intended to support beginners in Hindustani music. This has a large collection of Hindustani bandishes, with lyrics, notation, audio and information on raga, tala and laya.
% * BiLSTM
% * The network is pre-trained to continue the notes, and from here you get the embeddings.
% * Then you cluster the embeddings using K-means
% * 14 ragas

% \citep{dighe-swara_raga}
% * heuristic-based svara histogram from mixtures
% * nothing more really

% \citep{vasudevan-hybrid_raga}
% * combining a classification and a clustering model
% * Peak magnitude and phase values from short-time FFT windows (like a rudimentary pitch extraction)
% * no separation
% * 8 raagas

% \citep{bhattacharjee-raga_transition}
% * hand-annotated svara notes
% * 10 Hindustani ragas
% * Note sequences manually transcribed from 1 hour audio of musical pieces of 2 different performances
% * Transition Probability Matrix, Euclidean Distance and correlation coefficients. Note sequences manually transcribed from 1 hour audio of musical pieces of 2

 
% \citep{madhusudhan-deepsrgm}
% * Separate vocals using \citep{drossos2-mad_twinnet}, it's mono, and not optimized for music.
% * Praat API, optimized for speech and not robust to background noise
% * concatenated LSTM, trained on 10 and 40 ragas

% \citep{clayton-raga_classification}
% * Multimodal approach, combining visual and audio modalities
% * Use the four-stem spleeter
% * Use the Python API of PRAAT
% * 9 raagas

% \citep{vishnu-lstm_raga}
% * 5 different raagas
% * LSTM network
% * MFCC features
% * No separation

% \citep{john-carnatic_classification}
% * Spleeter 5-stems + noise removal
% * Praat API for pitch extraction
% * CNN
% * 5 raagas



% \subsection{Heuristic and Rule-Based Approaches}

% Early endeavors in raga recognition predominantly relied on heuristic methods, leveraging domain knowledge to model musical structures. For instance, \citet{pandey-tansen} developed the TANSEN system, which employed Hidden Markov Models (HMMs) combined with string matching techniques to identify rāgas based on note sequences derived from two heuristics: the Hill Peak and Note Duration heuristics. This approach focused on transcribing audio into note sequences without explicit pitch extraction or source separation.

% Similarly, \citet{shetty-raga} introduced a rule-based system that matched scales using autocorrelation for pitch detection. The system was designed for monophonic inputs, limiting its applicability to polyphonic compositions. \citet{dighe-swara_raga} proposed a heuristic classifier that utilized svara histograms extracted directly from audio mixtures, bypassing explicit pitch extraction and source separation.

% \citet{bhattacharjee-raga_transition} constructed transition probability matrices from manually annotated note sequences to model rāga transitions, emphasizing the importance of note sequences in rāga identification. \citet{vasudevan-hybrid_raga} combined classification and clustering models using peak magnitude and phase information from short-time Fourier transform (STFT) windows, serving as a rudimentary pitch extraction method without employing source separation.

% \subsection{Machine Learning-Based Approaches}

% With the advent of machine learning, researchers began to harness statistical models for raga recognition, often incorporating pitch-based features. \citet{chordia-raga} utilized pitch class distributions derived from Harmonic Product Spectrum (HPS) analysis to classify rāgas, focusing on tonal content without source separation.

% \citet{koduri-survey} analyzed note segmentation and pitch class histograms using Rao's pitch extraction method, providing insights into the tonal framework of rāgas. In a subsequent study, \citet{koduri-raga} combined YIN and Melodia algorithms for pitch extraction, applying k-nearest neighbors (k-NN) classification on pitch histograms to identify rāgas.

% \citet{kumar-raga} employed Support Vector Machines (SVMs) trained on pitch class profiles and n-gram histograms extracted using the Melodia algorithm, demonstrating improved classification performance. \citet{gulati-mining} developed a classifier ensemble leveraging tonal features from melodies extracted via Melodia, enhancing recognition accuracy.

% \citet{gulati-phrase_raga} applied HMM and Conditional Random Fields (CRFs) on phrase-level features, incorporating unvoiced region removal to refine pitch contours obtained from Melodia. \citet{gulati-time_delayed} introduced time-delayed melody surfaces (TDMS) as features, using Melodia for pitch extraction and demonstrating effectiveness in both Hindustani and Carnatic rāgas.

% \citet{belle-raga} extracted pitch contours using Rao's method, manually validating them to ensure accuracy, and utilized pitch class distributions for rāga classification. \citet{ross-raga_similarity} trained a BiLSTM network on symbolic note sequences to learn embeddings, subsequently clustering them using k-means to identify rāgas.

% \citet{jayanthi-raga} focused on spectral features such as spectral centroid, bandwidth, roll-off, zero-crossing rate, and MFCCs, employing a BiLSTM-XGBoost model without explicit pitch extraction or source separation. \citet{vishnu-lstm_raga} utilized MFCC features in conjunction with an LSTM model to classify rāgas, omitting pitch extraction and source separation steps. \citet{basu-raga} proposed a two-stage deep learning model based on MFCC features, achieving notable performance without relying on pitch extraction or source separation.

% \citet{chowdhuri-raga} implemented CNN and RNN models on chromagrams, indirectly capturing pitch information without direct pitch extraction or source separation. \citet{anand-raga} applied CNNs to chromogram features derived from Melodia-extracted pitch contours, focusing on predominant pitch values without source separation.

% \citet{shah-raga_recognition} combined CNN and LSTM architectures on spectrograms, employing Spleeter for source separation to isolate vocals, though the impact of separation on performance was not clearly established. \citet{shikarpur-raga_classification} analyzed mode switching in Hindustani music using pitch-based features extracted via PRAAT, incorporating four-stem Spleeter for source separation.

% \citet{madhusudhan-deepsrgm} trained an LSTM model on features from svara sequences, utilizing PRAAT for pitch extraction and Drossos et al.'s method for source separation, despite its limitations in music processing. \citet{clayton-raga_classification} adopted a multimodal approach, combining visual and audio features extracted via PRAAT, and applied four-stem Spleeter for source separation.

% \citet{john-carnatic_classification} employed a CNN on pitch features extracted using PRAAT, integrating five-stem Spleeter and denoising techniques for enhanced performance.